%=============================================================================
% Role 2: Solver and Numerical Methods
%
% LaTeX FRAGMENT -- no \documentclass, no preamble, no \begin{document}.
% \input this into the master document. Assumes the master preamble defines
% amsthm-style environments: proposition, lemma, remark, proof (plus amsmath),
% as used by the sibling parts in this directory.
% Notation matches Role 1: capacities K_R, K_F, K_M; costs C_R, C_F, C_M;
% prices/parameters R, c, beta, p; residual demand x with CDF F(x;K_R);
% density f(x;K_R); shift ratio g(x) = F_{K_R}(x)/f(x).
%=============================================================================

\label{part:role2}

Role 2 built and validated a numerical solver for the joint capacity problem,
independently re-deriving the paper's closed forms rather than transcribing
them, and using the solver as a lie-detector on the theory. Everything below is
grounded in code that was run: where a claim was \emph{proved} by hand it is
stated as a Proposition or Lemma; where it was only \emph{verified numerically}
that is said plainly, with the size of the check.

Throughout, the utility's expected profit is the single unified integral (paper
eq.~5, which coincides with eq.~10 in the large-FCL regime)
\begin{equation}
\label{eq:r2-profit}
\begin{aligned}
\Pi(K_R,K_F,K_M) =\;& R\!\int_{K_F-K_M}^{K_F}\! F(x;K_R)\,dx
   \;-\; c\!\int_{-K_M}^{K_F-K_M}\!\bigl(1-F(x;K_R)\bigr)\,dx \\[2pt]
   &-\;(\beta+p)\!\int_{K_F}^{\bar D}\!\bigl(1-F(x;K_R)\bigr)\,dx
   \;-\;\bigl(C_F K_F + C_R K_R + C_M K_M\bigr),
\end{aligned}
\end{equation}
taken against the \emph{extended} CDF $F(\cdot;K_R)$, defined to equal $0$ for
$x\le -K_R$ and $1$ for $x\ge \bar D$. The FCL dispatch rule is
$d(x)=\min\{\max(K_F-x,0),\,K_M\}$ (eq.~3). Carbon is proportional to expected
fossil output,
\begin{equation}
\label{eq:r2-carbon}
\mathrm{Carbon}(K_R,K_F,K_M) \;=\; c\!\int_{-K_M}^{K_F-K_M}\!\bigl(1-F(x;K_R)\bigr)\,dx,
\end{equation}
and the blackout probability is $1-F(K_F;K_R)$. The extension convention is what
lets one profit functional and one FOC system cover both the large-FCL regime
($K_M>K_R$, Section~4) and the small-FCL regime ($K_M\le K_R$, Appendix~G) with
no case split: outside the true support $F$ is simply clipped, and every
integral runs against that clipped $F$.

%-----------------------------------------------------------------------------
\section{Uniform closed forms and the no-finite-optimum finding}
\label{sec:r2-week1-uniform}

Under the uniform residual-demand family $F(x;K_R)=(x+K_R)/(\bar D+K_R)$ on
$[-K_R,\bar D]$ (Section~4.3, eqs~18--26), every optimality condition solves in
closed form. Write $Q=\bar D+K_R$.

\begin{proposition}[Uniform closed forms]
\label{prop:r2-uniform-cf}
For the uniform family the stationary capacities are:
\begin{align}
\text{no-FCL (eqs~20,21):}\quad
 & K_R^{o}=\bar D\!\left(\sqrt{\tfrac{c}{\,2C_R-C_F^2/(\beta+p-c)\,}}-1\right),
 && K_F^{o}=\bar D-\frac{C_F}{\beta+p-c}\,Q,\\[4pt]
\text{with-FCL at fixed $K_R$ (eqs~18,19$\equiv$22,23):}\quad
 & K_F^{*}=\frac{\beta+p-c-C_F-C_M}{\beta+p-R}\,Q-K_R,
 && K_M^{*}=K_F^{*}-\frac{C_M}{R-c}\,Q+K_R.
\end{align}
The associated expected fossil output (hence carbon, after multiplying by $c$)
is, with $a=K_F-K_M$,
\begin{align}
\text{no-FCL (eq.~25$\times c$):}\quad
 & \mathrm{Carbon}^{o}=c\,\frac{K_F^{o}\bigl(\bar D-K_F^{o}/2\bigr)}{Q},\\[4pt]
\text{with-FCL (eq.~26$\times c$, large-FCL):}\quad
 & \mathrm{Carbon}^{*}=c\!\left[(K_M-K_R)+\frac{\bar D\,(a+K_R)-\tfrac12(a^2-K_R^2)}{Q}\right].
\end{align}
\end{proposition}

Each formula was re-derived from the FOCs by hand-integrating
$F(x;K_R)=(x+K_R)/Q$ against the relevant limits, then cross-checked against the
general quadrature solver to machine precision. Three corrections to the OCR'd
source came out of that process.

\begin{remark}[Corrections found in re-derivation]
\label{rem:r2-corrections}
\emph{(a)} Equations~25 and 26 as printed are ``carbon$/c$'' -- expected fossil
\emph{output} in MWh, not emissions; multiplying the printed brackets by $c$
reproduces the quadrature carbon exactly. \emph{(b)} The with-FCL carbon in the
large-FCL regime is \emph{not} the small-FCL formula $K_F(\bar
D+K_M-K_F/2)/Q$ (Appendix~G eq.~76). When $K_M>K_R$ the range $[-K_M,K_F-K_M]$
sticks out past $-K_R$ by $(K_M-K_R)$, and on that stretch $1-F\equiv1$, giving
the flat ``$(K_M-K_R)$'' term above; this was cross-checked three ways
(quadrature, a from-scratch integral, and a literal re-typing of printed eq.~26)
to $8$ decimals. \emph{(c)} The knife-edge bracket of eq.~24 equals exactly
\emph{twice} a direct profit gradient (see \S\ref{sec:r2-week1-gradient}); the
factor of $2$ is a benign grouping difference that does not affect its sign.
\end{remark}

The central Week-1 finding is negative and structural.

\begin{proposition}[No interior optimum in the runaway sub-regime]
\label{prop:r2-no-finite-opt}
Under the uniform family (and its power-law generalization
$F=((x+K_R)/Q)^{\alpha}$), whenever $C_R<C_R^{\dagger}$ (equivalently the knife-edge
$\Delta_{24}>0$), the with-FCL profit \eqref{eq:r2-profit} has no finite interior
maximum: it increases without bound as $K_R$ and $K_M$ grow together. Under
Assumption~1's inequality $R<c+C_F+C_M$, the \emph{net coefficient of $K_R$} in the
closed form $K_F^{*}=\varphi_2 Q-K_R=\varphi_2\bar D-(1-\varphi_2)K_R$ is
$-(1-\varphi_2)<0$, so the interior stationary curve $K_F^{*}$ crosses zero at some
finite $K_R$: the positive case ``$K_R^{*}=\infty$'' of eq.~24 is a one-directional
first-order signal, not an attainable interior point. For $C_R>C_R^{\dagger}$
($\Delta_{24}<0$) the profit is instead maximized at the bounded corner $K_R^{*}=0$
(consistent with Role~4's boundedness threshold); the unboundedness is a property of
the cheap-renewable sub-regime, not universal.
\end{proposition}

The unboundedness was confirmed numerically (the joint profit keeps rising at
every increase of the optimization box, and the degeneracy extends to $K_F,K_M$
jointly, not just $K_R$). Because ``carbon at $K_R^{*}$'' is then ill-posed,
Role~2 instead swept the well-posed question: with a realistic cap on buildable
renewable capacity, does an optimized FCL program ever emit less carbon than the
(always bounded) no-FCL benchmark?

\begin{remark}[The 150-cell sweep]
\label{rem:r2-sweep}
Across the main grid of $150$ cells ($10$ values of $C_R$ $\times$ $5$ capacity
caps $\times$ $3$ distribution shapes), plus $60$ further cells over four
$(R,c,C_F,C_M)$ combinations -- $210$ cells in all -- profit and reliability
improved in \emph{every} cell, and carbon improved in \emph{zero} of them.
Proposition~1 parts (i)~profit-higher and (ii)~blackout-lower are
distribution-free and held whenever $C_M<\bar C_M=(R-c)\phi^{o}$, with
$\phi^{o}=F(K_F^{o};K_R^{o})$ the no-FCL reliability; part (iii) carbon-lower
never held in this search. This does not prove carbon improvement impossible in
general -- only that it does not appear anywhere in a wide, deliberate search
inside the model's economically valid domain.
\end{remark}

%-----------------------------------------------------------------------------
\section{The gradient identity and the analytic FOCs}
\label{sec:r2-week1-gradient}

The first-order conditions were obtained by hand-differentiating the unified
profit integral \eqref{eq:r2-profit} with Leibniz's rule, which makes them valid
in both regimes automatically.

\begin{lemma}[Analytic gradients]
\label{lem:r2-focs}
The profit gradient is
\begin{align}
\frac{\partial\Pi}{\partial K_M}
 &= (R-c)\,F(K_F-K_M;K_R) + c\,F(-K_M;K_R) - C_M,\\[2pt]
\frac{\partial\Pi}{\partial K_F}
 &= R\bigl[F(K_F)-F(K_F-K_M)\bigr] - c\bigl[1-F(K_F-K_M)\bigr]
    + (\beta+p)\bigl[1-F(K_F)\bigr] - C_F,\\[2pt]
\frac{\partial\Pi}{\partial K_R}
 &= R\!\int_{K_F-K_M}^{K_F}\!\!F_{K_R}\,dx + c\!\int_{-K_M}^{K_F-K_M}\!\!F_{K_R}\,dx
    + (\beta+p)\!\int_{K_F}^{\bar D}\!\!F_{K_R}\,dx - C_R,
\end{align}
where $F_{K_R}=\partial F/\partial K_R$. The first two are closed form
(differentiating an integral in its own limit just evaluates the integrand); only
$\partial\Pi/\partial K_R$ needs one quadrature, since $K_R$ reshapes $F$
everywhere rather than moving a limit.
\end{lemma}

At $K_M\ge K_R$ the term $F(-K_M;K_R)$ vanishes and $\partial\Pi/\partial K_M$
collapses to the paper's large-FCL eq.~12; at $K_M<K_R$ the same formula matches
Appendix~G eq.~61. The gradients were checked against central finite differences
of the quadrature profit for all three families (uniform, power
$\alpha\in\{0.5,2\}$, primitive), agreeing to the integration tolerance.

\begin{proposition}[Knife-edge identity, eq.~24]
\label{prop:r2-eq24}
Define the printed eq.~24 bracket
\begin{equation}
\Delta_{24}=(R-c)\!\left(\tfrac{R-C_M-c}{R-c}\right)^{2}+c
 +(\beta+p-R)\!\left(\tfrac{R-c-C_F-C_M}{\beta+p-R}\right)^{2}-2C_R,
\end{equation}
whose sign drives Section~4.3's trichotomy ($\Delta_{24}>0\Rightarrow
K_R^{*}=\infty$; $=0\Rightarrow$ any $K_R$ optimal; $<0\Rightarrow K_R^{*}=0$).
Evaluated along the curve $\bigl(K_F^{*}(K_R),K_M^{*}(K_R)\bigr)$, the direct
gradient $\partial\Pi/\partial K_R$ is \emph{constant in $K_R$} -- the actual
knife-edge content -- and $\Delta_{24}=2\,\partial\Pi/\partial K_R$ exactly.
\end{proposition}

Both the constancy and the factor-of-$2$ identity were verified numerically
across several parameter draws. The scale difference (the paper subtracts $2C_R$
where the raw gradient subtracts $C_R$ once) is a bookkeeping artifact of the
closed-form substitution and leaves the sign -- the only thing eq.~24 is used
for -- unchanged.

%-----------------------------------------------------------------------------
\section{The analytic Hessian and negative-definiteness certification}
\label{sec:r2-week2-hessian}

To ask whether a \emph{genuine} finite optimum exists, Role~2 built the
$3\times3$ Hessian analytically and certified negative-definiteness at candidate
stationary points. Since the problem is not globally concave
(Proposition~\ref{prop:r2-no-finite-opt}), the certificate is local: it confirms
a candidate is a true maximum rather than a saddle.

\begin{lemma}[Analytic Hessian]
\label{lem:r2-hessian}
Ordered $(K_R,K_F,K_M)$, the Hessian of $\Pi$ has entries (writing $f$ for the
density and $F_{K_R},F_{K_RK_R}$ for the first and second $K_R$-partials of $F$):
\begin{align*}
H_{MM}&=-(R-c)f(K_F-K_M)-c\,f(-K_M), &
H_{FF}&=\bigl[R-(\beta+p)\bigr]f(K_F)-(R-c)f(K_F-K_M),\\
H_{FM}&=H_{MF}=(R-c)f(K_F-K_M), &
H_{RM}&=H_{MR}=(R-c)F_{K_R}(K_F-K_M)+c\,F_{K_R}(-K_M),\\
H_{RF}&=H_{FR}=\bigl[R-(\beta+p)\bigr]F_{K_R}(K_F)-(R-c)F_{K_R}(K_F-K_M), &&
\end{align*}
\begin{equation*}
H_{RR}=R\!\int_{K_F-K_M}^{K_F}\!\!F_{K_RK_R}\,dx
 + c\!\int_{-K_M}^{K_F-K_M}\!\!F_{K_RK_R}\,dx
 + (\beta+p)\!\int_{K_F}^{\bar D}\!\!F_{K_RK_R}\,dx
 + (\text{boundary correction}).
\end{equation*}
The boundary correction adds, for each integral whose range strictly contains
$-K_R$, that term's coefficient times $F_{K_R}(-K_R;K_R)$ evaluated from inside
the support -- a moving-limit Leibniz term that is nonzero exactly in the
large-FCL regime $K_M>K_R$.
\end{lemma}

Each off-diagonal entry was derived twice (from two different gradients);
symmetry $H_{ij}=H_{ji}$ falls out rather than being assumed, which is itself an
algebra check. The whole Hessian was cross-checked against a central-difference
Jacobian of the analytic gradient (a single, not double, finite difference, so
far less noise), and the boundary term was caught by a failing finite-difference
test in precisely the $K_M>K_R$ regime. A quadrature subtlety mattered here: at
the support edge $-K_R$ the integrand has a breakpoint, and plain adaptive
\texttt{quad} silently returned values wrong by $\sim10^{-5}$ while reporting
near-zero error -- harmless for profit values but fatal for a second derivative.
Passing $-K_R$ as an explicit breakpoint (switching to a split-interval rule)
fixed it.

\begin{proposition}[Negative-definiteness certificate]
\label{prop:r2-negdef}
A symmetric $H$ is negative definite iff both criteria hold: \emph{(Sylvester)}
the signed leading principal minors $(-1)^{k}\det H_{[1:k,\,1:k]}>0$ for
$k=1,2,3$; and \emph{(eigenvalues)} every eigenvalue is $<0$. The two are
mathematically equivalent, so requiring both to agree turns any disagreement into
a bug detector (e.g.\ an asymmetric $H$ from an algebra slip).
\end{proposition}

\begin{proposition}[A finite interior optimum exists off the power family]
\label{prop:r2-finite-opt}
For the truncated-lognormal primitive family -- residual demand $x=D-K_R S$ with
$D\sim\mathrm{Lognormal}(\mu,\sigma)$ truncated to $(0,\bar D]$ and
$S\sim\mathrm{Beta}(2,2)$, which is \emph{not} a rescaling of one fixed shape in
$K_R$ -- the with-FCL problem does admit a genuine finite interior optimum. It is
certified by three independent checks, all of which must pass: \emph{(1)}
bound-independence (the solution stops moving as the optimization box is
enlarged, ruling out runaway); \emph{(2)} the analytic gradient vanishes; and
\emph{(3)} the Hessian is negative definite by
Proposition~\ref{prop:r2-negdef}.
\end{proposition}

Swapping the distribution family was necessary but not sufficient: $C_R$ still
had to be tuned away from the runaway regime, exactly as under the uniform
family. At any certified optimum Role~2 evaluated Role~1's endpoint condition in
terms of the \emph{shift ratio} $g(x)=F_{K_R}(x;K_R)/f(x;K_R)$, namely
\begin{equation}
\label{eq:r2-shiftratio}
g(K_F^{*}-K_M^{*}) \;\ge\; g(K_F^{*}),
\end{equation}
Role~1's sufficient condition for the FCL program to raise renewable capacity
($K_R^{*}>K_R^{o}$): the \emph{lower} dispatch endpoint $a=K_F^{*}-K_M^{*}$
carries the larger shift ratio, which holds automatically wherever $g$ is
non-increasing (regular families). \textup{(}This is the corrected direction;
an earlier draft carried the reversed inequality $g(K_F^{*})\ge g(a)$, fixed
across all roles to match Role~1's capacity-channel proof.\textup{)}
The ratio $g$ needs the density $f$ and the first
$K_R$-derivative $F_{K_R}$, both added analytically for every family (closed-form
for uniform/power/primitive; for the lognormal primitive via a closed-form
$f_D'(d)=-(f_D(d)/d)\bigl[1+(\ln d-\mu)/\sigma^{2}\bigr]$).

%-----------------------------------------------------------------------------
\section{Classification of critical points}
\label{sec:r2-week2-classify}

Sweeping $(C_R,\sigma)$ over the lognormal-primitive family, each cell was sorted
into one of five mutually exclusive outcomes by the three checks above.

\begin{remark}[Finite-optimum taxonomy]
\label{rem:r2-taxonomy}
The categories are: \emph{runaway} -- no finite optimum at all (the Week-1
failure mode, which persists here for small $C_R$, confirming it is a property of
the cost structure, not of the distribution); \emph{corner-zero} -- collapse to a
zero-capacity corner when $C_R$ is too high; \emph{optimizer-stalled} -- a
non-vanishing gradient at the largest box (a numerical, not mathematical, issue,
reported separately so it is never conflated with the two real categories);
\emph{saddle} -- a genuine stationary point ($\nabla\Pi\approx0$) whose Hessian is
not negative definite; and \emph{genuine-optimum} -- finite, interior, and
regular, the only category where the carbon question is well posed. Among all
genuine optima found in the sweep, carbon never fell relative to no-FCL, so the
Week-1 finding survives the move to a well-posed optimum.
\end{remark}

One family sits outside this taxonomy: for the power law with $\alpha<1$, once
$K_M>K_R$ the $K_R$-direction second derivative is not even finite ($F_{K_R}$
diverges at the support edge), so ``the Hessian is not negative definite'' is too
weak a description -- the Hessian does not exist. That case is documented and
tested separately rather than swept, since a nonexistent matrix cannot be
evaluated.

%-----------------------------------------------------------------------------
\section{The dispatch LP, greedy optimality, and the carbon contract}
\label{sec:r2-week3-lp}

Week~3 asked a sharper question: holding capacities fixed, can an
\emph{optimally designed} state-contingent dispatch contract lower carbon where
the greedy rule cannot? The admissible contract class is any measurable
$d:\mathbb{R}\to[0,K_M]$ with $d(x)=0$ for $x>K_F$ (consumer priority, since
$\beta+p>R$). Per state, fossil output is $z(x)=\min\{K_F,\max(x+d(x),0)\}$,
mining revenue is $R\,d(x)$, and blackout is $\max(x-K_F,0)$.

Discretizing the state $x_i$ with probability weights $w_i$ and introducing $z_i$
as a free variable pinned by $z_i\ge0$, $z_i\ge x_i+d_i$, $z_i\le K_F$ linearizes
the piecewise-linear $z$, so both dispatch problems are exact linear programs
(solved via HiGHS). At any minimizing objective the optimizer drives $z_i$ to its
tightest bound, recovering $z_i=\min\{K_F,\max(x_i+d_i,0)\}$ exactly.

\begin{lemma}[Firm-optimal dispatch is greedy]
\label{lem:r2-lemma1}
The firm-optimal contract maximizes $\sum_i w_i\bigl(R\,d_i-c\,z_i\bigr)$ over all
admissible $d(\cdot)$ subject to the consumer-priority cap $x_i+d_i\le K_F$ (mining
is served only into the headroom below the fossil ceiling; served mining $d_i$ never
displaces consumer supply). Under that cap $z_i=\max(0,x_i+d_i)$ and the objective is
increasing in $d_i$, so its unique solution is the greedy rule
$d^{*}(x)=\min\{\max(K_F-x,0),\,K_M\}$ of eq.~3: for $R>c$ it maximizes served $d_i$,
giving $d_i=K_M$ in the full-dispatch region $x_i+K_M\le K_F$ and $d_i=K_F-x_i$ in the
gap-filling zone where the cap binds.
\end{lemma}

The LP solution reproduced the greedy schedule to a relative carbon error below
$10^{-4}$ in every test, numerically confirming Lemma~\ref{lem:r2-lemma1}. The
carbon-optimal contract solves the mirror-image program,
\begin{equation}
\label{eq:r2-carbonlp}
\min_{d(\cdot)}\;\sum_i w_i\,c\,z_i
\qquad\text{s.t.}\qquad
\sum_i w_i\bigl(R\,d_i-c\,z_i\bigr)\ \ge\ \pi_{\text{floor}},
\end{equation}
with the floor set to the firm-optimal dispatch profit, so both solve the same
constrained problem from opposite directions. The carbon-optimal schedule
declines to dispatch in states where the marginal fossil draw is high, trading
revenue for lower emissions.

\begin{proposition}[The optimal contract does not flip the carbon sign]
\label{prop:r2-noflip}
At fixed profit-maximizing capacities, giving the FCL program the best possible
(carbon-minimizing) dispatch schedule still does not make its carbon fall below
the no-FCL benchmark. Across the swept grid the carbon-optimal contract produced
zero sign flips ($0$ of $20$ cells in the keystone fast sweep, and no flips in
any sweep run). Not only does the myopic greedy rule fail to flip carbon -- the
LP-optimal dispatch fails too.
\end{proposition}

For the full question, capacities and dispatch are optimized jointly as a bilevel
program (an outer derivative-free capacity search wrapping the exact inner LP),
minimizing carbon subject to the firm-optimal profit floor; the same no-flip
conclusion holds.

%-----------------------------------------------------------------------------
\section{LP duality certification}
\label{sec:r2-week3-duality}

Unlike the non-convex capacity NLP, the dispatch LP has a global optimum
certifiable exactly by duality, with no curvature condition, because a
finite-dimensional LP's KKT conditions are necessary and sufficient.

\begin{proposition}[Dispatch-LP optimality certificate]
\label{prop:r2-lpcert}
Each dispatch solution is certified by three quantities extracted from the
solver: \emph{strong duality} (the normalized primal--dual gap
$|\Pi_{\text{primal}}-\Pi_{\text{dual}}|/\max(1,|\Pi_{\text{primal}}|)$ below
$10^{-8}$); \emph{complementary slackness} (for every constraint, either it is
tight or its dual is zero, so $\max_i|\,\text{slack}_i\cdot\text{dual}_i|\approx0$);
and the \emph{shadow price} on the profit-floor constraint of
\eqref{eq:r2-carbonlp} -- the extra carbon that must be borne per unit of profit
floor, i.e.\ the marginal carbon cost of insisting on firm-optimal profit.
\end{proposition}

The certificate also reports which constraints bind at the optimum -- how many
states run fossil at capacity ($z_i=K_F$), how many dispatch at the cap
($d_i=K_M$), and how many sit at zero dispatch -- giving a structural read on the
solution beyond a single objective value. Every certified dispatch solution in
the Week-3 sweeps passed the duality-gap and slackness tests to machine
precision, and the profit-floor shadow price is the natural hand-off to Role~3's
comparative statics: it is exactly the local exchange rate between carbon and
profit at the contract optimum.
